
\documentclass[11pt]{article}

\usepackage[a4paper,margin=1in]{geometry}
\usepackage{amsmath,amssymb,amsthm,bm}
\usepackage{siunitx}
\usepackage{booktabs}
\usepackage{hyperref}
\usepackage{xcolor}
\usepackage{enumitem}
\usepackage{mathtools}

\hypersetup{
  colorlinks=true,
  linkcolor=blue!55!black,
  citecolor=blue!55!black,
  urlcolor=blue!55!black
}

% ---------------------------------------------------------------------
% SST / finite-core notation macros. These are kept explicit and local.
% ---------------------------------------------------------------------
\newcommand{\swirlarrow}{%
  \mathchoice{\mkern-2mu\scriptstyle\boldsymbol{\circlearrowleft}}%
             {\mkern-2mu\scriptstyle\boldsymbol{\circlearrowleft}}%
             {\mkern-2mu\scriptscriptstyle\boldsymbol{\circlearrowleft}}%
             {\mkern-2mu\scriptscriptstyle\boldsymbol{\circlearrowleft}}%
}
\newcommand{\vswirl}{\mathbf{v}_{\!\boldsymbol{\circlearrowleft}}}
\newcommand{\vnorm}{\lVert\mathbf{v}_{\!\boldsymbol{\circlearrowleft}}\rVert}
\newcommand{\rhof}{\rho_{\!f}}
\newcommand{\rhoE}{\rho_{\!E}}
\newcommand{\rhom}{\rho_{\!m}}
\newcommand{\rc}{r_c}
\newcommand{\LK}{\mathcal L_K}
\newcommand{\Ezero}{\mathcal E_0}
\newcommand{\Eeff}{\mathcal E_{\rm eff}}
\newcommand{\Xisph}{\Xi_{\rm sph}}
\newcommand{\alphacell}{\alpha_{\rm cell}}
\newcommand{\EpNLS}{E_p^{\rm NLS}}
\newcommand{\kappaNLS}{\kappa_{\rm NLS}^{\rm shell}}
\newcommand{\phigolden}{\varphi}

\newtheorem{theorem}{Theorem}
\newtheorem{lemma}{Lemma}
\newtheorem{definition}{Definition}
\newtheorem{proposition}{Proposition}
\newtheorem{remark}{Remark}

\sisetup{
  scientific-notation = true,
  exponent-product = \times,
  output-exponent-marker = \mathrm{e}
}

\title{From Calibration to Closure in a Compton-Core Finite-Filament Model:\\
Spherical and NLS Pressure-Cell Closure for the Fine-Structure Scale}

\author{Omar Iskandarani\\
\small Independent Researcher, Appingedam, The Netherlands}

\date{\today}

\begin{document}
\maketitle

\begin{abstract}
A previous Compton-core finite-filament scaling model showed that several
electron-scale quantities can be fixed algebraically once a dimensionless core
aspect scale is supplied.  This note develops the pressure-cell closure used by
the companion BEM/NLS calculation.  The strongest derived element is the
pressure-sector count \(N_p=4\), obtained from the spherical surface spectrum
and supported by nonlinear and global topology checks.  The pressure scale
\[
  E_p^{\rm NLS}
  =
  \frac{16\pi}{3}\mathcal L_K
  \left(1-\frac{11}{48\mathcal L_K^2}\right)
\]
is therefore a conditional finite-cell closure: the factor \(16\pi/3\) still
requires the sector-volume normalization \(4\pi/3\) per pressure sector per unit
ropelength, and the coefficient \(11/48\) requires the unweighted shell
normalization \(w_\perp=1\).  The latter is not yet derived from the full GP/NLS
core dynamics; the supplied relaxed-core proxy returns \(w_\perp\simeq2.076\).
For \(\mathcal L_K=16.371637\), the conditional scale is
\(E_p^{\rm NLS}=274.0748756568\).  This paper should therefore be read as a
structured pressure-cell/NLS-shell closure and gate isolation, not as a
stand-alone first-principles derivation of the electromagnetic fine-structure
constant.
\end{abstract}

\section{Purpose and status}

The objective is to separate a numerical calibration from a mathematical closure.  Let
\(\Ezero\) denote a zeroth-order topological aspect scale and let \(\LK=L_K/D\) denote the
ropelength of a diameter-\(D\) ideal trefoil centreline.  The present paper proves the following
conditional result:

\begin{quote}
Given the finite-core Biot--Savart self-energy, the circulation relation
\(\Gamma_0=2\pi \rc\vnorm\), the spherical pressure-cell functional of
Sec.~\ref{sec:spherical}, and the aspect-renormalization rule of Sec.~\ref{sec:aspect},
the numerical fine-structure scale is fixed by \(\Ezero\) and \(\LK\).
\end{quote}

The proof is conditional because \(\Ezero\) itself is not derived here from primitive field
equations.  Nor is the far-field electromagnetic \(1/r\) law derived.  Those are promoted to
explicit open theorems in Sec.~\ref{sec:fundamental}.  This keeps the paper falsifiable:
the model succeeds as a closure theorem if the stated assumptions are accepted, and it becomes
a fundamental derivation only if those assumptions are later obtained from a unique dynamical
principle.



The uniqueness of this finite-cell reduction is treated in the standalone
companion Paper~IV.  In the present paper the same assumptions are used as the
pressure-cell closure; Paper~IV proves that, within the stated minimal
\(SO(3)\)-covariant variational class, these assumptions uniquely select
\(N_p=4\), \(\sigma=11/3\), and \(\chi_R=2\).

\section{Derived-label gate and assumption ledger}
\label{sec:derived-label-gate}

The word ``derived'' is reserved here for quantities obtained from a stated
variational problem, boundary-value problem, or theorem without using the
empirical fine-structure constant, CODATA, \(m_e\), \(e\), \(R_\infty\), or
\(\hbar\) as numerical inputs.  Under this convention the local logarithmic
coefficient \(A_K\to1/(4\pi)\) is derived, but the pressure-cell coefficients
introduced below are not yet first-principles derived.  They are structured
closure hypotheses.

\begin{table}[h]
\centering
\caption{Status of the pressure-cell ingredients in this paper.}
\begin{tabular}{@{}llll@{}}
\toprule
Item & Expression & Current status & Requirement for ``derived'' label\\
\midrule
Local logarithmic coefficient &
\(A_K\to1/(4\pi)\) &
Theorem &
Already derived from the local Biot--Savart singularity\\
Zeroth pressure prefactor &
\(16\pi/3\) &
Minimal-shell closure hypothesis &
Derive or independently justify the \(\ell=0\oplus\ell=1\) pressure manifold\\
Finite-shell correction &
\(11/(48\mathcal L_K^2)\) &
Minimal-shell closure hypothesis &
Derive the unweighted \(3+2/3\) split from a controlled NLS/GP shell expansion\\
Cell radius closure &
\(R_{\rm cell}=2L_K\) &
Minimal-shell closure hypothesis &
Derive the self-dual inner/outer normalization or include \(\chi_R\)-sensitivity\\
\bottomrule
\end{tabular}
\label{tab:derived-label-ledger}
\end{table}

Thus this paper supplies a coherent analytical closure framework.  It does not
by itself prove that the electromagnetic fine-structure constant is derived
from first principles.  That stronger label requires deriving the entries
marked as closure hypotheses in Table~\ref{tab:derived-label-ledger}.


\subsection{Why the minimal pressure shell has four sectors}
\label{subsec:why-four-pressure-sectors}

The pressure manifold is taken to be the lowest \(SO(3)\)-invariant boundary
subspace capable of representing both uniform compression and displacement of
the cell pressure centre.  The pressure field on a spherical cell boundary
admits the harmonic expansion
\[
  p(\theta,\phi)
  =
  \sum_{\ell=0}^{\infty}\sum_{m=-\ell}^{\ell}
  p_{\ell m}Y_{\ell m}(\theta,\phi).
\]
Uniform compression is the monopole sector \(\ell=0\), with degeneracy
\(2\ell+1=1\).  The lowest pressure-gradient/displacement sector is
\(\ell=1\), with degeneracy \(3\).  Therefore the minimal pressure manifold is
\[
  \mathcal P_{\rm min}
  =
  \mathcal H_0\oplus\mathcal H_1,
  \qquad
  N_p=\dim\mathcal H_0+\dim\mathcal H_1=1+3=4.
\]
Modes \(\ell\ge2\) are quadrupolar and higher shape deformations.  In a
surface-stiffness action containing \(\int_{S^2}|\nabla_{S^2}p|^2d\Omega\),
they carry eigenvalues \(\ell(\ell+1)\) and enter as higher-order shape
corrections, not as zeroth-order pressure sectors.  The detailed coefficient
gates are summarized in Appendix~\ref{app:pressure-coefficient-targets}.
For a complete robustness claim, the companion numerical package should also
report an \(\ell\le2\) extension as a negative or sensitivity control.  If the
quadrupolar sector shifts the result at leading order, the \(\ell\le1\) cutoff
must be treated as a fit-relevant closure rather than a derived low-mode
truncation.

\subsection{Why the shell deficit is \(3+\langle\sin^2\theta\rangle\)}
\label{subsec:why-shell-deficit}

The finite-shell deficit is decomposed into a spherical volume second variation
and a transverse NLS-shell contribution,
\[
  \sigma=\sigma_{\rm vol}+\sigma_{\perp}.
\]
For a relative inward shell displacement \(\eta=a/R_{\rm cell}\), the spherical
volume Jacobian gives
\[
  \frac{V(R-a)}{V(R)}
  =
  (1-\eta)^3
  =
  1-3\eta+3\eta^2+O(\eta^3).
\]
After the first-order displacement has been absorbed into the equilibrium
radius, the second-order volume coefficient is
\[
  \sigma_{\rm vol}=3.
\]
The isotropic NLS/GP gradient term contributes through the transverse component
of the shell deformation.  For an isotropic shell orientation,
\[
  \sigma_\perp
  =
  \langle\sin^2\theta\rangle_{S^2}
  =
  \frac{\int_0^\pi \sin^2\theta\,\sin\theta\,d\theta}
       {\int_0^\pi \sin\theta\,d\theta}
  =
  \frac{2}{3}.
\]
Thus
\[
  \sigma=3+\frac{2}{3}=\frac{11}{3}.
\]
More generally one may write
\[
  \sigma=3+\frac{2}{3}w_\perp,
\]
where \(w_\perp=1\) is the isotropic NLS-gradient weighting used by the
minimal model.  This parameterization makes the remaining vulnerability
explicit: a first-principles derivation must obtain \(w_\perp=1\), not tune it
to the empirical fine-structure constant.
With \(\eta_K=1/(2\chi_R\mathcal L_K)\) and \(\chi_R=2\), this gives
\[
  \sigma\eta_K^2
  =
  \frac{11}{3}\frac{1}{16\mathcal L_K^2}
  =
  \frac{11}{48\mathcal L_K^2}.
\]

\subsection{Why \(A_\chi=\chi_R+N_p/\chi_R\) is the minimal radius action}
\label{subsec:why-radius-action}

Let \(\chi_R=R_{\rm cell}/L_K\).  The minimal inner/outer pressure balance has
two reciprocal costs.  The outer pressure-volume cost grows with the
dimensionless cell radius,
\[
  A_{\rm out}\propto \chi_R,
\]
whereas the inner crowding/compliance cost decreases with radius and is
proportional to the number of pressure sectors,
\[
  A_{\rm in}\propto \frac{N_p}{\chi_R}.
\]
With self-dual inner/outer normalization, the lowest-order reciprocal action is
\[
  A_\chi(\chi_R)
  =
  \chi_R+\frac{N_p}{\chi_R}.
\]
Stationarity gives
\[
  \frac{dA_\chi}{d\chi_R}
  =
  1-\frac{N_p}{\chi_R^2}=0,
  \qquad
  \chi_R=\sqrt{N_p}=2.
\]
A more general action \(a\chi_R+bN_p/\chi_R\) would give
\(\chi_R=\sqrt{bN_p/a}\).  Hence \(\chi_R=2\) is derived within the
self-dual inner/outer pressure normalization \(a=b\).  Without that
self-duality assumption, \(\chi_R=2\) remains a sensitivity parameter.


\subsection{Post-review gate-script updates}
\label{subsec:post-review-gate-updates}

Three gate calculations update the status of the pressure-cell coefficients.

First, the pressure-sector count is now supported by the fixed-volume surface
spectrum of a spherical cell.  For a perturbation
\(r(\Omega)=R[1+\epsilon f(\Omega)]\), the fixed-volume surface variation gives
\[
  \delta^2 A
  =
  \frac{R^2}{2}
  \int_{S^2}
  \left(|\nabla_{S^2}f|^2-2f^2\right)d\Omega .
\]
For \(f=Y_{\ell m}\),
\[
  k_\ell=\ell(\ell+1)-2=(\ell-1)(\ell+2).
\]
Thus \(\ell=1\) is the translation/centre-displacement zero-mode sector,
\(\ell\ge2\) are positive-stiffness shape modes, and \(\ell=0\) is the scalar
compression sector.  Hence
\[
  N_p=\dim\mathcal H_0+\dim\mathcal H_1=1+3=4
\]
is derived in the spherical finite-cell surface reduction.

This cutoff is also robust beyond the quadratic spectrum.  A perturbative
finite-core audit verifies that \(\ell\ge2\) remains separated from the retained
sector when \(\|\delta H\|_{\rm op}<k_2/2=2\), and the physical scale
\(\eta_K=1/(4\mathcal L_K)=1.52703116982\times10^{-2}\) is far below that
bound.  A nonlinear star-shaped solve through \(\ell_{\max}=6\) tested 45 real
\(\ell\ge2\) modes.  Across 300 random samples per amplitude and 20 local
minimizations from amplitude-one random starts, no competing leading
\(\ell\ge2\) pressure mode was found.  Global finite-perimeter cells are
protected by the isoperimetric inequality \(P(\Omega)^3\ge36\pi V(\Omega)^2\),
with equality only for a ball, so non-star-shaped, disconnected, or higher-genus
admissible cells cannot lower the fixed-volume area below the sphere.


Second, the GP core-profile second-variation audit solves the radial GP/NLS
vortex profile and confirms that the angular transverse factor is profile
independent under isotropic radial weights:
\[
  \left\langle\sin^2\theta\right\rangle_{S^2}=\frac23.
\]
The coefficient is therefore
\[
  \sigma=3+\frac23 w_\perp.
\]
For the unweighted isotropic shell normalization \(w_\perp=1\), one obtains
\[
  \sigma=\frac{11}{3},
  \qquad
  \frac{\sigma}{4\chi_R^2}=\frac{11}{48}
  \quad(\chi_R=2).
\]
Sensitivity runs with alternative profile weights do not generically reproduce
\(11/48\).  Hence the coefficient is derived within the unweighted isotropic
GP/NLS shell reduction, and the remaining microscopic gate is to derive
\(w_\perp=1\) from the full primitive reduction.

Third, the reciprocal radius action is now tied to the SST pressure identity.
With the canonical constants,
\[
  \frac12\rho_{\rm core}
  \lVert\mathbf v_{\!\boldsymbol{\circlearrowleft}}\rVert^2
  =
  \frac{F_{\rm swirl}^{\max}}{2\pi r_c^2}
  =
  2.329244600448\times10^{30}\ {\rm Pa}.
\]
Thus \(P_{\rm in}/P_{\rm out}=1\) in the leading Laplace-matched reciprocal
model, giving
\[
  A_\chi=\chi_R+\frac{N_p}{\chi_R},
  \qquad
  \chi_R=\sqrt{N_p}=2.
\]
The remaining robustness requirement is to exclude same-order nonreciprocal
terms such as \(\log\chi_R\), \(\chi_R^2\), or \(\chi_R^{-2}\).


\section{Claim-status correction: \(16\pi/3\) and \(11/48\)}
\label{sec:claim-status-16pi-11over48}

The surface-spectrum calculation derives \(N_p=4\), but this does not by itself
derive the numerical prefactor \(16\pi/3\).  The additional step is
\[
  E_p^{(0)}
  =
  N_p\left(\frac{4\pi}{3}\right)\mathcal L_K,
\]
which assumes a sector-volume normalization \(4\pi/3\) per pressure sector per
unit ropelength.  This is a compact and useful closure ansatz, but it is still a
load-bearing open coefficient until obtained from a variational pressure-cell
principle.

Similarly,
\[
  \sigma=3+\frac23w_\perp
\]
is the derived shell structure.  The value \(11/48\) follows only after setting
\(w_\perp=1\) and \(\chi_R=2\).  The current relaxed-core proxy gives
\(w_\perp\simeq2.076\), so the primitive GP/NLS derivation of \(w_\perp=1\) is
not established.  All precision claims that use \(11/48\) are therefore
conditional on the unweighted isotropic shell normalization.

\section{Finite-core filament energy}

Consider a smooth closed filament centreline \(X(s)\subset\mathbb R^3\), parametrized by
arclength \(s\in[0,L_K]\), with unit tangent \(T(s)=\partial_s X(s)\).  Let \(a\) be a small
core cutoff satisfying
\[
  a\ll \ell_{\rm curv},\qquad a\ll d_{\min},
\]
where \(\ell_{\rm curv}\) is the local curvature scale and \(d_{\min}\) is the minimum
nonlocal self-distance of the centreline.

The dimensionless cutoff Biot--Savart energy used in the numerical implementation is
\begin{equation}
  E_{\rm BS}^{\rm norm}(a)
  =
  \frac{1}{8\pi}
  \int_{0}^{L_K}\!\!\int_{0}^{L_K}
  \frac{T(s)\cdot T(s')}{\lVert X(s)-X(s')\rVert}
  \,\Theta\!\left(\lVert X(s)-X(s')\rVert-a\right)\,ds\,ds',
  \label{eq:dimensionless-bs}
\end{equation}
where \(\Theta\) is a cutoff indicator.  Its physical normalization is
\begin{equation}
  E_{\rm BS}(a)
  =
  \rhof\Gamma_0^2
  E_{\rm BS}^{\rm norm}(a).
\end{equation}
For a slender filament, the fitted local form is
\begin{equation}
  E_{\rm BS}(a)
  =
  \rhof\Gamma_0^2 L_K
  \left[
  A_K\log\!\left(\frac{L_K}{a}\right)+a_K
  \right]
  +
  o(1)
  \quad(a/L_K\to0).
  \label{eq:bs-log-form}
\end{equation}

\section{Local logarithmic coefficient}
\label{sec:AK}

\begin{theorem}[Universal local Biot--Savart logarithm]
For any \(C^2\) smooth slender closed filament, the leading local logarithmic coefficient in
Eq.~\eqref{eq:bs-log-form} is
\[
  A_K=\frac{1}{4\pi}
\]
provided the energy is normalized as in Eq.~\eqref{eq:dimensionless-bs}.
\end{theorem}

\begin{proof}
Fix \(s\) and write \(u=s'-s\).  In the near-diagonal region
\(a<|u|<\ell\), where \(a\ll \ell\ll\ell_{\rm curv}\), smoothness gives
\[
  X(s+u)-X(s)=uT(s)+O(u^2),
  \qquad
  T(s)\cdot T(s+u)=1+O(u^2).
\]
Therefore
\[
  \frac{T(s)\cdot T(s+u)}{\lVert X(s+u)-X(s)\rVert}
  =
  \frac{1}{|u|}+O(1).
\]
The local singular part of Eq.~\eqref{eq:dimensionless-bs} per unit arclength is then
\[
  \frac{1}{8\pi}
  \int_{a<|u|<\ell}\frac{du}{|u|}
  =
  \frac{1}{8\pi}
  \left[
  2\log\!\left(\frac{\ell}{a}\right)
  \right]
  =
  \frac{1}{4\pi}\log\!\left(\frac{\ell}{a}\right).
\]
Nonlocal contributions and curvature-dependent finite terms alter \(a_K\), not the coefficient
of \(\log(1/a)\).  Hence \(A_K=1/(4\pi)\).
\end{proof}

\begin{remark}
This theorem is local.  It does not assert that every finite-resolution numerical plateau is
exactly \(1/(4\pi)\); it asserts the slender-core asymptotic coefficient.  Knot geometry enters
the finite part \(a_K\), the ropelength \(\LK\), and nonlocal corrections.
\end{remark}

\section{Pressure balance and core closure}
\label{sec:pressure}

The normal pressure exerted by the exterior irrotational self-energy on a tube of radius \(a\)
is the negative derivative with respect to volume.  Since the tube area is \(2\pi aL_K\),
\begin{equation}
  P_{\rm irr}(a)
  =
  -\frac{1}{2\pi aL_K}\frac{\partial E_{\rm BS}}{\partial a}.
\end{equation}
Using Eq.~\eqref{eq:bs-log-form},
\begin{equation}
  P_{\rm irr}(a)
  =
  \frac{\rhof\Gamma_0^2A_K}{2\pi a^2}.
  \label{eq:Pirr}
\end{equation}
The core kinetic pressure scale is
\begin{equation}
  P_{\rm core}
  =
  \frac{1}{2}\rhof\vnorm^2.
  \label{eq:Pcore}
\end{equation}
The no-contact pressure balance \(P_{\rm irr}=P_{\rm core}\) gives
\begin{equation}
  \frac{\rhof\Gamma_0^2A_K}{2\pi a^2}
  =
  \frac{1}{2}\rhof\vnorm^2,
\end{equation}
hence
\begin{equation}
  a
  =
  \frac{\Gamma_0}{\vnorm}\sqrt{\frac{A_K}{\pi}}.
\end{equation}
With the circulation relation
\begin{equation}
  \Gamma_0=2\pi \rc\vnorm,
  \label{eq:gamma-relation}
\end{equation}
this becomes
\begin{equation}
  \boxed{\frac{a_{\rm nc}}{\rc}=\sqrt{4\pi A_K}.}
  \label{eq:anc}
\end{equation}
Therefore, in the theorem limit \(A_K=1/(4\pi)\),
\begin{equation}
  \boxed{a_{\rm nc}=\rc.}
\end{equation}

Dimensional consistency is immediate:
\[
  [P_{\rm irr}]
  =
  \frac{[\rhof][\Gamma_0]^2}{[a]^2}
  =
  \frac{\mathrm{kg\,m^{-3}}\;\mathrm{m^4\,s^{-2}}}{\mathrm{m^2}}
  =
  \mathrm{kg\,m^{-1}\,s^{-2}}
  =
  \mathrm{Pa}.
\]

\section{Spherical pressure-cell functional}
\label{sec:spherical}

The local calculation fixes the tube radius.  The next question is whether the exterior
spherical cell can generate the finite correction required for the fine-structure scale.

Let a trefoil tube of total length \(L_K\) and diameter \(D=2a\) be embedded in a spherical
coarse-graining cell of radius
\begin{equation}
  R_{\rm cell}=2L_K.
  \label{eq:Rcell}
\end{equation}
Then
\begin{equation}
  \eta_K
  =
  \frac{a}{R_{\rm cell}}
  =
  \frac{a}{2L_K}
  =
  \frac{1}{4(L_K/D)}
  =
  \frac{1}{4\LK}.
  \label{eq:eta}
\end{equation}

\subsection{Geometric pressure-transfer factor}

A small outward radial displacement of the spherical control surface by \(a\) gives the
surface-area Jacobian
\begin{equation}
  J_A(\eta_K)
  =
  \frac{4\pi(R_{\rm cell}+a)^2}{4\pi R_{\rm cell}^2}
  =
  (1+\eta_K)^2
  =
  1+2\eta_K+\eta_K^2.
  \label{eq:JA}
\end{equation}
The isotropic radial volume-response mode contributes one additional linear factor,
\begin{equation}
  J_V^{(r)}(\eta_K)=\eta_K.
  \label{eq:JV}
\end{equation}
The total spherical pressure-transfer multiplier is therefore
\begin{equation}
  \boxed{
  \Xisph
  =
  J_A+J_V^{(r)}
  =
  1+3\eta_K+\eta_K^2
  =
  1+\frac{3}{4\LK}+\frac{1}{16\LK^2}.
  }
  \label{eq:Xi}
\end{equation}

\subsection{Explicit variational form}

To make the above statement variational rather than merely kinematic, introduce an isotropic
cell-response coordinate \(q\).  Let \(K_{\rm sph}>0\) be a spherical compression modulus and
define
\begin{equation}
  \mathcal G_{\rm sph}(q;\eta_K)
  =
  \frac{K_{\rm sph}}{2}q^2
  -
  K_{\rm sph}\left[J_A(\eta_K)+J_V^{(r)}(\eta_K)\right]q.
  \label{eq:G-sph}
\end{equation}
Stationarity gives
\begin{equation}
  \frac{\partial\mathcal G_{\rm sph}}{\partial q}
  =
  K_{\rm sph}
  \left[
  q-\left(J_A+J_V^{(r)}\right)
  \right]
  =
  0,
\end{equation}
hence
\begin{equation}
  q_\star=\Xisph.
\end{equation}
Since
\begin{equation}
  \frac{\partial^2\mathcal G_{\rm sph}}{\partial q^2}=K_{\rm sph}>0,
\end{equation}
the stationary point is a strict minimum.

\begin{remark}
Equation~\eqref{eq:G-sph} is an explicit coarse-grained pressure-cell functional.  It does not
yet derive \(R_{\rm cell}=2L_K\) from primitive dynamics; Eq.~\eqref{eq:Rcell} is a closure
choice.  A fully fundamental derivation must replace Eq.~\eqref{eq:Rcell} by a stationarity
condition on \(R_{\rm cell}\).
\end{remark}


The derivation-target lemmas for the pressure prefactor, the finite-shell coefficient, and the radius closure are collected in Appendix~\ref{app:pressure-coefficient-targets}.  This appendix is part of the derived-label gate: it identifies the exact mode-count and shell-variation results still required to upgrade the pressure-cell scale from structured closure to first-principles derivation.

\section{NLS finite-shell pressure scale}
\label{sec:nls-shell}

The spherical response factor in Sec.~\ref{sec:spherical} fixes the finite-core transfer
geometry, but the companion finite-cell eigenvalue problem also requires a pressure
equilibrium scale and a stiffness.  The non-linear Schrödinger/Gross--Pitaevskii vortex-core
asymptotics provide a natural source for both quantities.  For a large vortex ring with core
radius \(r_c\), the standard asymptotic form is
\begin{equation}
  E_{\rm ring}(R)
  =
  \frac{1}{2}\rho\,\Gamma^2 R
  \left[
    \ln\!\left(\frac{8R}{r_c}\right)-A
  \right],
  \label{eq:nls-ring-energy}
\end{equation}
and the translational velocity has the form
\begin{equation}
  V_{\rm ring}(R)
  =
  \frac{\Gamma}{4\pi R}
  \left[
    \ln\!\left(\frac{8R}{r_c}\right)-B
  \right].
  \label{eq:nls-ring-velocity}
\end{equation}
Hamiltonian consistency \(V_{\rm ring}=dE_{\rm ring}/dP\), with
\(P(R)=\rho\Gamma\pi R^2\), gives
\begin{equation}
  B=A-1.
  \label{eq:nls-hamiltonian-constraint}
\end{equation}
The Golden NLS branch uses
\begin{equation}
  A=\varphi,\qquad B=\varphi^{-1},
  \qquad
  \varphi=\frac{1+\sqrt5}{2}.
  \label{eq:golden-branch}
\end{equation}

The zeroth-order spherical pressure scale associated with an ideal trefoil cell is
\begin{equation}
  E_p^{(0)}
  =
  \frac{16\pi}{3}\mathcal L_K.
  \label{eq:Ep0}
\end{equation}
Finite tube thickness contributes a second-order shell correction in the small parameter
\begin{equation}
  \eta_K=\frac{1}{4\mathcal L_K}.
  \label{eq:eta-nls}
\end{equation}
The NLS finite-shell pressure scale used in the companion BEM calculation is
\begin{equation}
  \boxed{
  \EpNLS
  =
  \frac{16\pi}{3}\mathcal L_K
  \left[
  1-\frac{11}{48\mathcal L_K^2}
  \right]
  =
  \frac{16\pi}{3}\mathcal L_K
  \left[
  1-\frac{11}{3}\eta_K^2
  \right].
  }
  \label{eq:EpNLS}
\end{equation}
The corresponding shell stiffness is
\begin{equation}
  \boxed{
  \kappaNLS
  =
  8\pi\mathcal L_K^2
  =
  \frac{\pi}{2\eta_K^2}.
  }
  \label{eq:kappaNLS}
\end{equation}
The coefficient \(11/48\) is not fitted to \(\alpha\), but it is conditional.
The post-review GP core-profile audit derives the decomposition
\(\sigma=3+(2/3)w_\perp\).  The unweighted isotropic shell normalization
\(w_\perp=1\) gives \(\sigma=11/3\) and therefore
\(11/(48\mathcal L_K^2)\) at \(\chi_R=2\).  However, the current relaxed-core
proxy gives \(w_\perp\simeq2.076\).  Thus a full primitive GP/SST core-profile
reduction must still explain why the shell normalization relevant to the
finite-cell pressure response is \(w_\perp=1\), or else the \(11/48\) correction
must be replaced.  The stiffness
\(\kappa_{\rm NLS}^{\rm shell}\) is the corresponding thin-shell stiffness; it is large
because \(\eta_K\ll1\).

The finite-cell pressure enthalpy used by the companion BEM eigenvalue calculation is
\begin{equation}
  \mathcal A_{\rm pressure}(E)
  =
  \kappaNLS
  \left[
    \frac{1}{3}\left(\frac{E}{\EpNLS}\right)^3
    -
    \ln\!\left(\frac{E}{\EpNLS}\right)
  \right].
  \label{eq:pressure-action-nls}
\end{equation}
It satisfies
\begin{equation}
  \frac{d\mathcal A_{\rm pressure}}{d\log E}
  =
  \kappaNLS
  \left[
    \left(\frac{E}{\EpNLS}\right)^3-1
  \right],
  \label{eq:pressure-derivative-nls}
\end{equation}
so the pressure part alone has a strict stationary point at \(E=\EpNLS\).  The BEM part in
the companion paper supplies the finite-cell compliance correction:
\begin{equation}
  \mathcal A_{\rm total}(E)
  =
  \mathcal A_{\rm BEM}(E)
  +
  \mathcal A_{\rm pressure}(E),
  \qquad
  \frac{d\mathcal A_{\rm total}}{d\log E}=0.
  \label{eq:total-action-nls}
\end{equation}
Thus the present paper supplies the analytical NLS-shell closure, while the companion
calculation determines the stationary eigenvalue \(E_\star\).

For the ideal trefoil value \(\mathcal L_K=16.371637\), Eq.~\eqref{eq:EpNLS} gives
\begin{align}
  E_p^{(0)} &= 274.309410808,\\
  \EpNLS &= 274.074875657,\\
  \kappaNLS &= 6736.341149.
\end{align}
These are dimensionless geometric/NLS closure quantities; no electron mass, Planck constant, charge,
Rydberg constant, or CODATA value is used to evaluate them.  Their use as first-principles inputs remains conditional on deriving the prefactor \(16\pi/3\), the shell coefficient \(11/48\), and the cell-radius closure.

\section{Aspect renormalization and fine-structure scale}
\label{sec:aspect}

The pressure-cell closure assumes that the zeroth-order aspect scale \(\Ezero\) receives a
small finite-core correction proportional to the local core coefficient and to the spherical
transfer response.  In the theorem limit \(A_K=1/(4\pi)\),
\begin{equation}
  \pi^2 A_K\to\frac{\pi}{4}.
\end{equation}
The effective aspect scale is therefore
\begin{equation}
  \boxed{
  \Eeff
  =
  \Ezero
  \left[
  1-
  \frac{\pi}{4\Ezero^2}\Xisph
  \right].
  }
  \label{eq:Eeff}
\end{equation}
The closure identification is
\begin{equation}
  \boxed{
  \alphacell=\frac{2}{\Eeff}.
  }
  \label{eq:alpha-cell}
\end{equation}
Combining Eqs.~\eqref{eq:Xi}, \eqref{eq:Eeff}, and \eqref{eq:alpha-cell} gives
\begin{equation}
  \boxed{
  \alphacell
  =
  \frac{2}{
  \Ezero
  \left[
  1-\dfrac{\pi}{4\Ezero^2}
  \left(
  1+\dfrac{3}{4\LK}+\dfrac{1}{16\LK^2}
  \right)
  \right]
  }.
  }
  \label{eq:main-alpha}
\end{equation}

\begin{remark}
Equation~\eqref{eq:main-alpha} is a theorem inside the specified closure model.  It is not yet
an unconditional derivation of the QED coupling because the far-field interaction law has not
been derived.  The necessary additional theorem is stated in Sec.~\ref{sec:fundamental}.
\end{remark}

\section{Numerical evaluation}
\label{sec:numerical-evaluation}

Use the ideal-trefoil ropelength value
\begin{equation}
  \mathcal L_K = 16.371637,
\end{equation}
taken from the ideal-knot Fourier data source.  Then
\begin{align}
  \eta_K &= \frac{1}{4\mathcal L_K}=0.01527031170,\\
  \Xi_{\rm sph} &= 1+3\eta_K+\eta_K^2=1.046044118,\\
  E_p^{(0)} &= \frac{16\pi}{3}\mathcal L_K=274.309410808,\\
  \EpNLS
  &=
  \frac{16\pi}{3}\mathcal L_K
  \left(1-\frac{11}{48\mathcal L_K^2}\right)
  =
  274.074875657,\\
  \kappaNLS &= 8\pi\mathcal L_K^2=6736.341149.
\end{align}

The companion BEM/NLS calculation uses \(\EpNLS\) and \(\kappaNLS\) in
Eq.~\eqref{eq:total-action-nls} and finds the stationary aspect eigenvalue \(E_\star\).  Once
that eigenvalue is supplied, the finite-core effective aspect is evaluated as
\begin{equation}
  E_{\rm eff}
  =
  E_\star
  \left[
  1-
  \frac{\pi}{4E_\star^2}
  \left(
  1+\frac{3}{4\mathcal L_K}+\frac{1}{16\mathcal L_K^2}
  \right)
  \right],
  \label{eq:Eeff-star}
\end{equation}
and the dimensionless fine-structure-scale estimate is
\begin{equation}
  \alpha_{\rm cell}=\frac{2}{E_{\rm eff}}.
  \label{eq:alpha-star}
\end{equation}

\begin{table}[h]
\centering
\caption{Analytical pressure-cell quantities supplied to the companion BEM/NLS finite-cell eigenvalue calculation.}
\begin{tabular}{@{}lll@{}}
\toprule
Quantity & Expression & Value \\
\midrule
Ideal trefoil ropelength & \(\mathcal L_K\) & \(16.371637\) \\
Shell parameter & \(\eta_K=1/(4\mathcal L_K)\) & \(0.01527031170\) \\
Spherical transfer & \(\Xi_{\rm sph}\) & \(1.046044118\) \\
Zeroth pressure scale & \(E_p^{(0)}=(16\pi/3)\mathcal L_K\) & \(274.309410808\) \\
NLS pressure scale & \(\EpNLS\) & \(274.074875657\) \\
NLS shell stiffness & \(\kappaNLS=8\pi\mathcal L_K^2\) & \(6736.341149\) \\
\bottomrule
\end{tabular}
\label{tab:nls-pressure}
\end{table}



\subsection{Phase--pressure normal-mode closure of the far-field Hessian}

The far-field stiffness gate has also been sharpened by the companion
normal-mode audit.  The exterior Hodge calculation gives
\[
  q_\phi=1,\qquad \Lambda_\phi=4\pi R\mathcal K_{\rm cell}.
\]
The interior chain now proceeds as follows.  The radial phase problem gives the
optimal profile \(n_{\rm eff}(r)\propto r^{-2}\).  The accessible-area
projection gives
\[
  M_{\rm acc}=M_\phi\frac{(R-a)^2}{R^2}.
\]
Finally, in the canonical linearized Madelung/GP normal-mode reduction,
\[
  H_k=\frac12p_k^2+\frac12\omega_k^2q_k^2
\]
has equal cycle-averaged phase/flow and pressure/compression budgets:
\[
  M_\phi=M_p=\frac{E_{\rm eff}}{2}.
\]
Therefore
\[
  \Lambda_\phi=\frac{E_{\rm eff}}{2R^2},
\]
and for \(R=1\),
\[
  \Lambda_\phi=\frac{E_{\rm eff}}{2}.
\]
This closes the far-field Hessian gate within the canonical normal-mode
reduction.  The remaining bridge assumption is that the \(E_{\rm eff}\) used in
the pressure-cell chain is the cycle-averaged quadratic energy of that same
interior canonical mode family.

\section{Reproducibility files for the coefficient gates}
\label{sec:gate-reproducibility-files}

The gate calculations referenced in this paper are stored in \texttt{../code/}
relative to the \texttt{Manuscripts} directory:
\begin{itemize}
\item \(N_p=4\): \texttt{derive\_pressure\_mode\_cutoff\_surface\_spectrum.py},
\texttt{test\_finite\_core\_nonspherical\_shape\_corrections.py}, and
\texttt{solve\_nonlinear\_shape\_stability.py};
\item \(11/48\) conditional on \(w_\perp=1\):
\texttt{derive\_gp\_core\_profile\_second\_variation.py};
\item \(w_\perp\) warning check:
\texttt{derive\_relaxed\_gp\_core\_hessian\_wperp.py};
\item \(\chi_R=2\): \texttt{derive\_pressure\_self\_duality\_from\_laplace\_matching.py};
\item \(q_\phi=1\) and exterior capacity:
\texttt{solve\_one\_cell\_hodge\_phase\_hessian.py};
\item \(\Lambda_\phi=E_{\rm eff}/2\) within the canonical phase--pressure
normal-mode reduction:
\texttt{derive\_phase\_pressure\_exchange\_self\_duality.py}.
\end{itemize}
The corresponding CSV and Markdown outputs are in the matching
\texttt{../code/outputs\_*} directories.

\section{What remains for a fundamental derivation}

A further independent selection principle is required for the trefoil itself.
The present paper uses \(K=3_1\) because that is the candidate finite-core
filament assigned to the electron sector in the model.  This assignment is not
derived here.  The falsifiable requirement is to run the same closure for the
unknot, the figure-eight knot \(4_1\), and other low-complexity knots and to
show that they do not reproduce the observed fine-structure scale under the
same coefficient rules.

\label{sec:fundamental}

The preceding sections provide a conditional mathematical closure.  To upgrade it to a
fundamental derivation, the following results must be obtained without using \(\alpha\),
\(R_\infty\), \(e\), or CODATA as input.


\subsection{Updated open gates after the post-review scripts}

The current open gates are narrower than in the original closure statement.
The pressure-sector count \(N_p=4\) is derived in the spherical surface
variation, and \(\chi_R=2\) is derived in the Laplace-matched reciprocal
pressure model.  What remains is:
\begin{enumerate}[label=(G\arabic*)]
\item derive the unweighted transverse normalization \(w_\perp=1\) from the
full GP/SST core-profile reduction;
\item extend the nonlinear shape-stability audit beyond star-shaped finite-perimeter cells only if the physical pressure-cell admissibility class is enlarged beyond reduced boundaries;
\item identify the \(E_{\rm eff}\) entering the pressure-cell chain with the
cycle-averaged quadratic energy of the same canonical interior mode family used
in the phase--pressure normal-mode reduction.
\end{enumerate}

\subsection{Open Theorem I: topological aspect eigenvalue}

The number
\begin{equation}
  \Ezero=274.074996
\end{equation}
must be obtained as an eigenvalue of a well-posed geometric or hydrodynamic variational
problem:
\begin{equation}
  \delta \mathcal F[X,\mathbf u,p,R_{\rm cell},a]=0,
  \qquad
  \nabla\cdot\mathbf u=0,
\end{equation}
with trefoil topology and circulation constraints.  Without this theorem, \(\Ezero\) remains
an external aspect input.

\subsection{Open Theorem II: spherical radius from stationarity}

The closure radius \(R_{\rm cell}=2L_K\) should be replaced by
\begin{equation}
  \frac{\partial \mathcal F}{\partial R_{\rm cell}}=0.
\end{equation}
If this stationarity condition yields \(R_{\rm cell}=2L_K\), then
\(\eta_K=1/(4\LK)\) becomes dynamical rather than chosen.

\subsection{Open Theorem III: far-field interaction law}

A fine-structure constant is not only a number.  In conventional electrodynamics it is the
coefficient of a long-range coupling.  The required far-field theorem is
\begin{equation}
  V_{\rm eff}(r)
  =
  \frac{\alphacell\hbar c}{r}
  +
  O(r^{-2})
  \qquad (r\to\infty).
  \label{eq:farfield}
\end{equation}
Equivalently, the model must yield
\begin{equation}
  \alphacell
  =
  \frac{g_{\rm eff}^2}{4\pi\hbar c}.
\end{equation}
Only after Eq.~\eqref{eq:farfield} is derived from the same variational structure can
Eq.~\eqref{eq:main-alpha} be called a fundamental derivation of the electromagnetic coupling.

\section{Falsification criteria}

The closure chain is falsified, or at least not fundamental, if any of the following fail:
\begin{enumerate}[label=(F\arabic*)]
\item The slender-core coefficient does not converge to \(A_K=1/(4\pi)\).
\item The ideal-trefoil ropelength used for \(\LK\) is not the relevant minimizing geometry.
\item A pressure-cell variational problem does not yield \(R_{\rm cell}=2L_K\) or an equivalent
      \(\eta_K=1/(4\LK)\).
\item The aspect-renormalization law Eq.~\eqref{eq:Eeff} cannot be derived from a more primitive
      pressure or action principle.
\item The far-field interaction is not of the \(1/r\) form in Eq.~\eqref{eq:farfield}.
\end{enumerate}

\section{Conclusion}


The standalone uniqueness proof in Paper~IV shows that the pressure manifold,
shell coefficient, and radius closure are not independent tunings once the
minimal finite-cell variational reduction is imposed.  Without Paper~IV they
should be read as closure hypotheses; with Paper~IV they are consequences of
that stated reduction.


The pressure-coefficient gates are stated explicitly in Appendix~\ref{app:pressure-coefficient-targets}; until those gates are closed, the NLS pressure scale should be called closure-derived rather than first-principles derived.


The mathematical route from calibration to pressure-cell closure is now separated into three
levels.  First, the local Biot--Savart singularity gives the slender-core coefficient
\(A_K\to1/(4\pi)\), and the pressure balance then fixes \(a_{\rm nc}=r_c\).  Second, the
spherical control-volume response gives
\[
  \Xi_{\rm sph}
  =
  1+\frac{3}{4\mathcal L_K}+\frac{1}{16\mathcal L_K^2}.
\]
Third, the NLS finite-shell closure supplies the pressure scale and stiffness
\[
  \EpNLS
  =
  \frac{16\pi}{3}\mathcal L_K
  \left(1-\frac{11}{48\mathcal L_K^2}\right),
  \qquad
  \kappaNLS=8\pi\mathcal L_K^2.
\]
For the ideal trefoil these evaluate to
\[
  \EpNLS=274.074875657,
  \qquad
  \kappaNLS=6736.341149.
\]
These quantities are computed from the ropelength geometry and the stated NLS-shell closure hypotheses, not from
\(\alpha\), \(m_e\), \(\hbar\), \(e\), or \(R_\infty\).  They should therefore be labelled as closure-derived rather than first-principles derived until the coefficients in Table~\ref{tab:derived-label-ledger} are obtained from an independent variational calculation.  The stationary aspect eigenvalue
\(E_\star\) is obtained in the companion BEM/NLS finite-cell calculation.  The final
identification with the electromagnetic fine-structure constant further requires the
far-field interaction theorem
\[
  V(r)=\frac{\alpha_{\rm cell}\hbar c}{r}+O(r^{-2}).
\]

\appendix

\section{Dimensional checks}

The three key energy terms have dimensions
\begin{align}
  [\rhof\Gamma_0^2L_K]
  &=
  \mathrm{kg\,m^{-3}}\,
  \mathrm{m^4\,s^{-2}}\,
  \mathrm{m}
  =
  \mathrm{J},\\
  \left[\frac{1}{2}\pi\rhof\vnorm^2a^2L_K\right]
  &=
  \mathrm{kg\,m^{-3}}\,
  \mathrm{m^2\,s^{-2}}\,
  \mathrm{m^2}\,
  \mathrm{m}
  =
  \mathrm{J},\\
  [P_{\rm irr}]
  &=
  \mathrm{Pa}.
\end{align}
The quantities \(A_K\), \(\LK\), \(\eta_K\), \(\Xisph\), \(\Ezero\), \(\Eeff\), and
\(\alphacell\) are dimensionless.

\section{Small-parameter expansion}

Let
\[
  \epsilon=\frac{1}{\Ezero}.
\]
Then Eq.~\eqref{eq:Eeff} is
\[
  \Eeff
  =
  \Ezero
  \left[
  1-\frac{\pi}{4}\Xisph\epsilon^2
  \right].
\]
Therefore
\[
  \alphacell
  =
  \frac{2}{\Ezero}
  \left[
  1+\frac{\pi}{4}\Xisph\epsilon^2
  +O(\epsilon^4)
  \right].
\]
Since \(\epsilon\simeq3.65\times10^{-3}\), the correction is naturally of order
\(10^{-5}\), matching the observed difference between the zeroth-order
\(2/\Ezero\) estimate and the refined closure.


\section{Derivation-target lemmas for pressure-cell coefficients}
\label{app:pressure-coefficient-targets}

This appendix records the exact gates that must be closed before the
pressure-cell scale can be labelled first-principles derived.  The statements
below are intentionally split into algebraic implications, which are proved,
and physical inputs, which remain derivation targets.

\subsection{Lemma A1: spherical pressure-volume factor}

Let \(\mathcal L_K=L_K/D\) be the ideal-trefoil ropelength.  Suppose the
zeroth-order pressure-cell action decomposes into \(N_p\) equivalent spherical
pressure sectors, each contributing the normalized spherical volume factor
\(4\pi/3\) per unit ropelength.  Then
\begin{equation}
  E_p^{(0)}
  =
  N_p\,\frac{4\pi}{3}\,\mathcal L_K .
  \label{eq:app-Ep0-general-Np}
\end{equation}
Consequently,
\begin{equation}
  E_p^{(0)}
  =
  \frac{16\pi}{3}\mathcal L_K
  \quad\Longleftrightarrow\quad
  N_p=4.
  \label{eq:app-Ep0-Np4}
\end{equation}

\paragraph{Proof.}
Equating Eq.~\eqref{eq:app-Ep0-general-Np} to the target coefficient gives
\[
  N_p\frac{4\pi}{3}=\frac{16\pi}{3},
\]
hence \(N_p=4\). \(\square\)

\paragraph{Derived-label gate.}
The coefficient \(16\pi/3\) is derived only after \(N_p=4\) is derived from a
pressure-cell variational principle or an independent boundary-mode count.  In
the present paper \(N_p=4\) is therefore a structured closure target.

\subsection{Lemma A2: NLS finite-shell second variation}

Let
\begin{equation}
  \eta_K
  =
  \frac{a}{R_{\rm cell}}
  =
  \frac{1}{2\chi_R\mathcal L_K},
  \qquad
  R_{\rm cell}=\chi_R L_K,
  \label{eq:app-eta-chi-general}
\end{equation}
where \(D=2a\).  Assume a second-order shell response
\begin{equation}
  E_p^{\rm shell}
  =
  E_p^{(0)}
  \left[
  1-\sigma\eta_K^2+O(\eta_K^3)
  \right].
  \label{eq:app-Ep-shell-general}
\end{equation}
Then
\begin{equation}
  E_p^{\rm shell}
  =
  E_p^{(0)}
  \left[
  1-\frac{\sigma}{4\chi_R^2\mathcal L_K^2}
  +O(\mathcal L_K^{-3})
  \right].
  \label{eq:app-Ep-shell-c2}
\end{equation}
For \(\chi_R=2\), the coefficient \(11/(48\mathcal L_K^2)\) is obtained iff
\begin{equation}
  \sigma=\frac{11}{3}.
  \label{eq:app-sigma-required}
\end{equation}

\paragraph{Proof.}
Substitution of Eq.~\eqref{eq:app-eta-chi-general} into
Eq.~\eqref{eq:app-Ep-shell-general} gives
\[
  \sigma\eta_K^2
  =
  \frac{\sigma}{4\chi_R^2\mathcal L_K^2}.
\]
At \(\chi_R=2\), this is \(\sigma/(16\mathcal L_K^2)\).  Matching
\[
  \frac{\sigma}{16}
  =
  \frac{11}{48}
\]
gives \(\sigma=11/3\). \(\square\)

A concrete NLS-shell target is the decomposition
\begin{equation}
  \sigma
  =
  3+\frac{2}{3}
  =
  \frac{11}{3}.
  \label{eq:app-sigma-decomposition}
\end{equation}
The first contribution is assigned to the second variation of the spherical
volume/pressure response.  The second contribution is assigned to isotropic
transverse shell averaging, e.g.
\[
  \langle \sin^2\theta\rangle_{S^2}=\frac{2}{3}.
\]
A future GP/NLS core-profile calculation must produce this decomposition
without reference to the empirical fine-structure constant.

\subsection{Minimal radius-closure target}

The current pressure-cell construction uses
\[
  R_{\rm cell}=2L_K.
\]
Introduce \(\chi_R=R_{\rm cell}/L_K\) and the minimal balancing action
\begin{equation}
  A_\chi(\chi_R)
  =
  \chi_R+\frac{\lambda_\chi}{\chi_R}.
  \label{eq:app-Achi-minimal}
\end{equation}
Then
\begin{equation}
  \frac{dA_\chi}{d\chi_R}
  =
  1-\frac{\lambda_\chi}{\chi_R^2}=0
  \quad\Longrightarrow\quad
  \chi_R=\sqrt{\lambda_\chi}.
\end{equation}
Thus
\begin{equation}
  \chi_R=2
  \quad\Longleftrightarrow\quad
  \lambda_\chi=4.
  \label{eq:app-chi-target}
\end{equation}
The open task is to derive \(\lambda_\chi=4\) from the inner/outer
pressure-cell balance, rather than imposing \(\chi_R=2\).


\subsection{Model-justification notes}

The preceding gates are consequences of the minimal shell model, not arbitrary
number insertion.  The use of \(\mathcal H_0\oplus\mathcal H_1\) is forced by
the requirement that the boundary pressure represent both scalar compression
and vector displacement while retaining only the lowest \(SO(3)\)-closed
subspace.  The coefficient \(3+2/3\) is the sum of a volume-Jacobian
second-variation coefficient and the isotropic transverse NLS average.  The
radius action \(A_\chi=\chi_R+N_p/\chi_R\) is the reciprocal self-dual
inner/outer pressure balance.  These statements are sufficient for a
``derived within the minimal finite-cell shell model'' label.  A stronger
first-principles label requires deriving the minimal shell model itself from
the underlying NLS/SST field equations.

\subsection{Summary of Paper II derived-label gates}

The pressure-cell route may be summarized as
\[
  \mathcal L_K
  \xrightarrow[\text{target: }N_p=4]{}
  \frac{16\pi}{3}\mathcal L_K
  \xrightarrow[\text{targets: }\sigma=11/3,\ \chi_R=2]{}
  E_p^{\rm NLS}.
\]
Thus Paper II supplies a structured closure framework.  It becomes a
first-principles pressure-scale derivation only after \(N_p=4\),
\(\sigma=11/3\), and \(\chi_R=2\) are obtained from an independent
pressure-shell variational calculation.


\begin{thebibliography}{99}

\bibitem{Helmholtz1858}
H. von Helmholtz, ``Über Integrale der hydrodynamischen Gleichungen, welche den Wirbelbewegungen entsprechen,''
\emph{Journal für die reine und angewandte Mathematik} \textbf{55}, 25--55 (1858).
English translation: ``On Integrals of the Hydrodynamical Equations, which Express Vortex-motion,''
\emph{Philosophical Magazine} \textbf{33}, 485--512 (1867).

\bibitem{Kelvin1867}
W. Thomson (Lord Kelvin), ``On Vortex Atoms,''
\emph{Proceedings of the Royal Society of Edinburgh} \textbf{6}, 94--105 (1867).

\bibitem{Saffman1992}
P. G. Saffman, \emph{Vortex Dynamics}, Cambridge University Press, Cambridge (1992).
ISBN: 978-0521477390.

\bibitem{Batchelor1967}
G. K. Batchelor, \emph{An Introduction to Fluid Dynamics}, Cambridge University Press, Cambridge (1967).
ISBN: 978-0521663960.

\bibitem{Cantarella2002}
J. Cantarella, R. B. Kusner, and J. M. Sullivan,
``On the minimum ropelength of knots and links,''
\emph{Inventiones Mathematicae} \textbf{150}, 257--286 (2002).
DOI: 10.1007/s00222-002-0234-y.
Preprint: arXiv:math/0103224.

\bibitem{Cantarella2011}
J. Cantarella, J. H. G. Fu, R. B. Kusner, and J. M. Sullivan,
``Ropelength criticality,''
\emph{Geometry \& Topology} \textbf{18}, 1973--2043 (2014).
DOI: 10.2140/gt.2014.18.1973.
Preprint: arXiv:1102.3234.

\bibitem{KnotAtlasIdeal}
D. Bar-Natan and contributors,
``The Knot Atlas: Ideal knot data file \texttt{Ideal.txt.gz},''
available at \url{https://katlas.org/images/d/d2/Ideal.txt.gz}.
Accessed 2026.

\bibitem{Jackson1999}
J. D. Jackson, \emph{Classical Electrodynamics}, 3rd ed.,
Wiley, New York (1999). ISBN: 978-0471309321.

\bibitem{PeskinSchroeder1995}
M. E. Peskin and D. V. Schroeder,
\emph{An Introduction to Quantum Field Theory},
Addison-Wesley, Reading, MA (1995). ISBN: 978-0201503975.

\bibitem{Pitaevskii1961}
L. P. Pitaevskii,
``Vortex Lines in an Imperfect Bose Gas,''
\emph{Soviet Physics JETP} \textbf{13}, 451--454 (1961).

\bibitem{RobertsGrant1971}
P. H. Roberts and J. Grant,
``Motions in a Bose Condensate. I. The Structure of the Large Circular Vortex,''
\emph{Journal of Physics A: General Physics} \textbf{4}, 55--72 (1971).
DOI: 10.1088/0305-4470/4/1/306.

\bibitem{PethickSmith2008}
C. J. Pethick and H. Smith,
\emph{Bose--Einstein Condensation in Dilute Gases}, 2nd ed.,
Cambridge University Press, Cambridge (2008).
DOI: 10.1017/CBO9780511802850.

\bibitem{MohrCODATA2022}
P. J. Mohr, D. B. Newell, B. N. Taylor, and E. Tiesinga,
``CODATA Recommended Values of the Fundamental Physical Constants: 2022,''
arXiv:2409.03787 (2024).
Permalink: \url{https://arxiv.org/abs/2409.03787}.

\bibitem{Osserman1978}
R. Osserman,
``The isoperimetric inequality,''
\emph{Bulletin of the American Mathematical Society} \textbf{84},
1182--1238 (1978).
DOI: 10.1090/S0002-9904-1978-14553-4.

\bibitem{Maggi2012}
F. Maggi,
\emph{Sets of Finite Perimeter and Geometric Variational Problems},
Cambridge University Press, Cambridge (2012).
DOI: 10.1017/CBO9781139108133.

\bibitem{Madelung1927}
E. Madelung,
``Quantentheorie in hydrodynamischer Form,''
\emph{Zeitschrift f\"ur Physik} \textbf{40}, 322--326 (1927).
DOI: 10.1007/BF01400372.

\bibitem{Gross1961}
E. P. Gross,
``Structure of a Quantized Vortex in Boson Systems,''
\emph{Nuovo Cimento} \textbf{20}, 454--477 (1961).
DOI: 10.1007/BF02731494.

\bibitem{Bogoliubov1947}
N. N. Bogoliubov,
``On the Theory of Superfluidity,''
\emph{Journal of Physics USSR} \textbf{11}, 23--32 (1947).

\end{thebibliography}

\end{document}
